% AI-written, AI cross-reviewed; no human sign-off; not a claim that RH is proved
% RH Notes blueprint source excerpts for Z1.3.04
% Verbatim excerpts in the notation of the parent blueprint. The labelled editions and source ranges appear below.
% See the unit YAML for the current proof status. Earlier-round passages retain their original wording.

% Start of atlas edition annotation. The following block is outside the original excerpt.
\paragraph{Status of named inputs in this reading.}
\noindent\nolinkurl{UpperTheory.V.Z.Kato.Joint.kato_global_unit_law_of_smoothed_classes} --- in-repo constant (\path{Development/AutomorphicArithmetic/AutomorphicVarieties/ModularityArgument/KatoQuadraticCharacters/FrobeniusHodgeTateLaws/KatoGlobalUnitLawOfSmoothedClasses.lean:8--10}); not compared.\par
% End of atlas edition annotation.

\section*{Z1.3.04. One global system with local constraints}

% Parent source SHA-256: 1af64922b98ed2310744411b1f18b25cc426c97220646321471836e1a3a409d6
% blueprint/ch_P.tex lines 13247 through 13255; r6:Z1.3.04
\subsection{Z1.3.04. Simultaneous smoothing of the global and local data}\label{r6:Z1.3.04}
Given the class theorem, the smoothed law and the admissible projection statement, start with the \emph{same} ordinary source throughout. The eigenfunctional construction chooses $N,\chi,\theta,\alpha$ and the two smoothing integers $c,d$. Put $c_w(x)=\iota^{-1}(\widehat\theta(x)/2)$. The Fourier input gives a nonzero coefficient and the integrality bound. Adjoin the finitely many $c_w(x)$ to a coefficient ring already containing the entries of $W$; apply admissibility there, choosing $\mathrm{pr},c_f,\lambda,\kappa_A$ together. The marked matrices supply the required bottom rows. The class theorem gives one family $Y$, its specialization $y$ and one power $\ell^m$ for both. The modified diamond character agrees on every residue unit, so the same functional and projection satisfy the required period identity. Apply the local law to this $Y$ at $d_0=\ell+1$; it gives $\kappa_5\ne0$ and the compatible periods $v_m$.
\par{\footnotesize \nolinkurl{UpperTheory.V.Z.Kato.Joint.kato_global_unit_law_of_smoothed_classes} (\nolinkurl{Development/AutomorphicArithmetic/AutomorphicVarieties/ModularityArgument/KatoQuadraticCharacters/FrobeniusHodgeTateLaws/KatoGlobalUnitLawOfSmoothedClasses.lean}\allowbreak:8--10). \nolinkurl{UpperTheory.V.Z.Kato.Joint.glue2} (\nolinkurl{Development/AutomorphicArithmetic/AutomorphicVarieties/ModularityArgument/KatoQuadraticCharacters/FrobeniusHodgeTateLaws/AdmissibleHeckeAlgebraOfGenerated.lean}\allowbreak:17--86).\par}

Let $b_i,t_i$ be the original four smoothing coefficients and frequencies. Append the four pairs $(-d_0b_i,d_0t_i)$. All eight frequencies remain prime to $\ell$. In the twisted moment at exponent $e$ their sum is $(1-d_0^{e+1})S_4$: the character of $d_0$ is one since $d_0\equiv1\bmod\ell$. The original four-term estimate gives $S_4\ne0$, and $d_0>1$ gives the other factor. The tame averaging identity contributes $\varphi(N)$, so the nonzero comparison scalar before multiplication by $\ell^m$ is $\kappa_5\kappa_A\varphi(N)$. Rescaling uses the same exponent on the global cocycle and on the smoothed second coordinate. The source inertia identity and determinant formula give $\nu_D|_I=1$; the Frobenius input supplies a lift with unit value. These are the data used to obtain the pin over an enlarged integral coefficient ring. In parallel, linearity of the smoothing and $d_0\equiv1\pmod\ell$ preserve the three system and specialization identities. Extending coefficients sends these identities to that same ring. This yields one tuple supporting all four assertions.
\par{\footnotesize \nolinkurl{UpperTheory.V.Z.Kato.Joint.glue2} (\nolinkurl{Development/AutomorphicArithmetic/AutomorphicVarieties/ModularityArgument/KatoQuadraticCharacters/FrobeniusHodgeTateLaws/AdmissibleHeckeAlgebraOfGenerated.lean}\allowbreak:87--134). \nolinkurl{UpperTheory.V.Z.Kato.Joint.glue2} (\nolinkurl{Development/AutomorphicArithmetic/AutomorphicVarieties/ModularityArgument/KatoQuadraticCharacters/FrobeniusHodgeTateLaws/AdmissibleHeckeAlgebraOfGenerated.lean}\allowbreak:153--186). \nolinkurl{UpperTheory.V.Z.Kato.Joint.glue2} (\nolinkurl{Development/AutomorphicArithmetic/AutomorphicVarieties/ModularityArgument/KatoQuadraticCharacters/FrobeniusHodgeTateLaws/AdmissibleHeckeAlgebraOfGenerated.lean}\allowbreak:187--230).\par}

\emph{Partial. [Proved in Lean, not reproduced:]} the ordinary eigenfunctional giving the initial $\theta,\alpha,c,d$ with its measure and smoothing clauses, and the construction converting the measured character moments, the explicit law, and the nonzero eight-term moment into a pin for this smoothed family. Those two inputs are supplied internally to the theorem, in addition to its three stated hypotheses.
\par{\footnotesize \nolinkurl{UpperTheory.V.Z.Kato.Joint.kato_global_unit_law_of_smoothed_classes} (\nolinkurl{Development/AutomorphicArithmetic/AutomorphicVarieties/ModularityArgument/KatoQuadraticCharacters/FrobeniusHodgeTateLaws/KatoGlobalUnitLawOfSmoothedClasses.lean}\allowbreak:8--10). \nolinkurl{UpperTheory.V.Z.Kato.Joint.glue2} (\nolinkurl{Development/AutomorphicArithmetic/AutomorphicVarieties/ModularityArgument/KatoQuadraticCharacters/FrobeniusHodgeTateLaws/AdmissibleHeckeAlgebraOfGenerated.lean}\allowbreak:21--43). \nolinkurl{UpperTheory.V.Z.Kato.Joint.glue2} (\nolinkurl{Development/AutomorphicArithmetic/AutomorphicVarieties/ModularityArgument/KatoQuadraticCharacters/FrobeniusHodgeTateLaws/AdmissibleHeckeAlgebraOfGenerated.lean}\allowbreak:197--230).\par}

% Parent source SHA-256: 1fdc6ea9aac6cf10316f85b3253b33ae98289a76ff5aa2e79daaaf11e057ce18
% blueprint/ch_P_r7.tex lines 752 through 774; r7:Z1.3.04
\subsection{Z1.3.04. One eigenfunctional and one measured tuple}\label{r7:Z1.3.04}
\textbf{B06.} Construct the initial ordinary eigenfunctional at one
auxiliary level with its smoothers, character and ordinary root.
There is a residue obstruction to Kato 12.6(2) at that level: at
\(e=1\) the prescribed pair \(n_1=1,n_2=0\) requires
\(d^2\ne d^2(\theta\chi)(d)\), hence \(d\not\equiv1\pmod N\).
The form's original level may divide \(N\); the calculation concerns the
auxiliary level. Section~\ref{r6:kato-lit} records the differently
indexed symbols in 12.6(1).
\nolinkurl{GlobalTheory.EisM.CuspidalUnitaryTwistsClaim},
\path{Statements/CuspidalUnitaryTwists/CuspidalUnitaryTwistsClaim.lean:13--25};
\nolinkurl{GlobalTheory.T.ordinaryEisensteinEigenfunctional_exists},
\path{FoundationProofs/OrdinaryEisensteinEigenfunctionalExistsStatement.lean:11--12}.

\textbf{B07.} Obtain \(\mathrm{Pin7}\) for the measured tuple, with one
nonzero scalar \(\kappa_0/\kappa\), an integral series, the two signed
character formulas and the stipulated lower bound. The source takes the
series and \(\kappa\) from the measure input, embeds both coefficient
rings in one, and increases the common starting level for both signs.
Kato 12.5(1)'s character values do not identify this measure or furnish
its uniform lower bound on this tuple. See Section~\ref{r6:Z1.3.04},
\nolinkurl{UpperTheory.V.Z.Kato.Joint.pinOf},
\path{FrobeniusHodgeTateLaws/LawMono.lean:72--97}.

